\section{Running the Adventure}
\label{sec:runningAdventure}

To run this adventure you will need \handbook{} containing all rules and guidelines on how to run \cosmereRPG{}. This book contains only adventure-specific details and suggestions on how to run it.

\quoteBox{
	You can read (or paraphrase) a text in such a box to your players. Alternatively, you can describe the scene using your own words.
}

Whenever a characters name appears in \enemy{bold}, it means there is corresponding stat block available. You can use it to get more info about the character's attributes and abilities.

Items appearing in \loot{bold italic} have special rules you can find at the end of this book.

\reasonBox{
	Boxes like this will help you understand the reason behind some of the design choices. Those boxes have no impact on the game, but they may help you modify some elements of the adventure if you so choose.
}

\ptrBox{
	When a plot might allow for certain Radiant Spren to be attracted by player's actions, that will be indicated by such a box. However, use your own judgment whenever a player does something not accounted by this book.
}

\corruptedBox{
	Some events or actions are relevant only for Enlightened Spren - the ones corrupted by \npcSjaAnat{}'s touch. Information about those will be displayed in such a box.
}

\spoilerBox{
	On many occasions, this book will reference events or characters from \bookSA{} series. This box will warn you of spoilers, but also it will provide original source if you wished to read more details about what's referenced from official books.
}

If you are using this book in electronic version, it has also several quality-of-life improvements to make your life easier when you navigate this book:

Whenever a name of a chapter or section appears (eg. "\nameref{sec:runningAdventure}"), you can click on it to quickly go to the chapter in question.

All characters' names (eg. \npcSjaAnat{}) are links that lead to \hyperref[personae]{Dramatis Person\ae} chapter, where all the NPCs are listed.

Items (eg. \lootDustcaster{}) and adversaries (eg. \enemyCultist{}) are also links that will lead you to the page with their unique rules. 



\subsection{Adjusting Difficulty}

This book assumes your players will form a party of 4 characters starting at 1st level. However, it is perfectly fine to have a different size of your group and/or to start playing at a different level (eg. because the characters already shared another adventure together). Also, it is possible that challenges presented in this book are not aligning with strong suits of the party.

If you see that the balance is a bit off, feel free to adjust the difficulty. You can change a number or tier of the enemies, modify endeavors' requirements or weave some other changes into the story.



\subsection{Handling Canon}
\label{ssec:runningAdventureCanon}

Despite treating original source with respect, this adventure is not canon. While playing the game your primary focus should be fun. For everyone, "fun" means something different - either following official lore to the letter, creating your own touching story, or just hitting bad guys with your sword. Make sure you know what is fun for you and your group and make it your primary objective. If that means modifying this adventure or straying away from the canon - go for it!

You can find more information on that topic in “Canon and Continuity” section of \handbook{}.




\subsection{Playing with Singers}

It is possible for players to play as Singers. However, they must be aware that for the first chapters of the story, they would need to hide their form and pretend to be parshmen not to raise too many questions from the Kholinar people.

Starting with \nameref{chap:chapter3}, after the arrival of Everstorm, Parshmen are out of the picture. Hiding Singer's form is no longer required, but the Kholinar soldiers are actively fighting the newly-turned Singers and that may introduce new challenges to Singer characters. On the other hand, they can find new allays, possibly opening new paths to reach the party's goals.



\subsection{Playing with Radiants}

The adventure encourages players to play as Radiants. Please work with with your players to know what Knight Order would they wish to join. 

Unless the player explicitly states different wishes, assume they will bond Enlightened Spren touched by \npcSjaAnat{}. Pay attention to the PC's actions that might attract a Spren. Even if no opportunity occurs to introduce them, they will be able to bond their Spren at the end of \nameref{chap:chapter2}.



\subsection{Unique Rules}

This adventure introduces possibility to bond Enlightened Spren touched by \npcSjaAnat{}. For those, you can use any rules you prefer, but it is assumed you will be following the ones provided in \hyperref[chap:corruptedRules]{Appendix \ref{chap:corruptedRules}}.

\subsubsection{Suggested Expedition Events}

If you are using Enlightened Radiant rules from \hyperref[chap:corruptedRules]{Appendix \ref{chap:corruptedRules}}, you will be expected to expedite the Radiant's progression when they meet specific Unmade or their influence.

Each time an Enlightened Radiant is exposed to an Unmade related to their Order, they can immediately develop a talent from their Enlightened talent tree. Each Radiant can be rewarded like that up to 3 times.

Here is a list of suggested moments to do that so. You can adhere to those suggestion or modify something to your needs:

\begin{itemize}
	\item \textbf{Enlightened Windrunner (\npcThrill{}):} \\
			\fullref{sec:bloodyStreets}	\\
			\fullref{sec:enemyReborn} 	\\
			\nameref{chap:chapter9}
			
	\item \textbf{Enlightened Skybreaker (\npcYeligNar{}):} \\
			\fullref{sec:sjaanatIntroduction}	\\
			\fullref{sec:closedPalace}	\\
			\nameref{chap:chapter9}
			
	\item \textbf{Enlightened Dustbringer (\npcDustmother{}):} \\
			\fullref{ssec:monasteryReliquary}	\\
			\fullref{sec:dustcastingSneaking}	\\
			\fullref{sec:dustcasting}
			
	\item \textbf{Enlightened Edgedancer (\npcMidnightMother{}):} \\
			\fullref{sec:visionsEnterShadesmarMidnight}	\\
			\fullref{ssec:stoppingWrongnessMidnightEssence}	\\
			\fullref{ssec:middnightMotherFight}
			
	\item \textbf{Enlightened Truthwatcher (\npcMoelach{}):} \\
			\fullref{sec:bloodyStreets}	\\
			\fullref{ssec:pathToKharbranthTrouble}		\\
			\fullref{ssec:conclaveHospitalPatient}
			
	\item \textbf{Enlightened Lightweaver (\npcSjaAnat{}):} \\
			\fullref{sec:sjaanatIntroduction}	\\
			\fullref{sec:dustcastingBriefing}	\\
			\nameref{chap:chapter9}
			
	\item \textbf{Enlightened Elsecaller (\npcBlackFisher{}):} \\
			\fullref{ssec:monasteryReliquary}	\\
			\fullref{ssec:pathToLighthouseGuide}		\\
			\fullref{sec:toThePhysicalRealm}
			
	\item \textbf{Enlightened Willshaper (\npcHeartRevel{}):} \\
			\fullref{sec:monasteryFeast}	\\
			\fullref{sec:revelry}	\\
			\fullref{sec:heartOfTheRevel}
			
	\item \textbf{Enlightened Stoneward (\npcMishram{}):} \\
			\fullref{sec:visionsEnterShadesmar}	\\
			\fullref{ssec:visionsExitShadesmar}	\\
			\nameref{chap:chapter10}
\end{itemize}



\subsubsection{\abilitySprenName{}}
\label{ssec:ruleSpren}

While traversing the Cognitive Realm, the characters will encounter many spren. In the \bookSA{} it is stated that they cannot be killed like other living creatures can. 

To reflect this, they gain this ability:

\abilitySpren{}

For the combat purposes, after reaching 0 Health, the creature with this feature is gone just like any other enemy. 

However, that does not mean it was killed. Consider them running away, or yielding the fight. Either way, the party doesn't need to worry about them anymore.

Please note, that there exist very specific weapons capable of killing the spren, either in Cognitive or Physical Realm. In those cases, this feature will not save them from their doom.

\subsubsection{\abilityFutureKnowledgeName{}}
\label{ssec:ruleDiagram}

In a moment of supreme intelligence, \npcTaravangian{} wrote "The Diagram" - set of cryptic prophecies treating about the fate of Roshar and future events.

Any character allied with him gains access to this knowledge, as a result gaining the feature below:

\abilityFutureKnowledge{}

If the party manages to ally the king, they will also gain access to the Diagram and its knowledge. However, working afar may make the access to the prophecies more difficult, so the PCs must use the knowledge sparingly.

This ability works the same as Enlightened Truthwatcher's \talentTruthwatcherAlterFortuneName{} talent, except it requires an access to The Diagram to refresh.

There might be also alternative sources of knowledge about the future: Death Rattles, fortune reading, and so on...
